% 第5章图：弦图（边12,13,23,24,34,35,45）的极大团A={1,2,3},B={2,3,4},C={3,4,5}；交集权 AB=2, BC=2, AC=1。
% 最大权生成树 A-B-C 总权4；B-A-C 总权3，含4的团B,C被A隔开。
\begin{tikzpicture}[x=1mm,y=1mm,>=Stealth,
  v/.style={circle,draw=FigureStroke,fill=FigureBlueFill,line width=.7pt,minimum size=5.4mm,inner sep=0pt},
  e/.style={FigureStroke,line width=.7pt},
  cl/.style={rectangle,rounded corners=1.8pt,draw=FigureStroke,fill=FigureBlueFill,line width=.7pt,inner sep=1.5pt,minimum height=6mm},
  badcl/.style={cl,draw=FigureRed,fill=FigureRedFill},
  w/.style={font=\figurefont\footnotesize,text=FigureInk,fill=white,inner sep=.8pt},
  sc/.style={font=\figurefont\footnotesize,text=FigureInk},
  note/.style={font=\FigureNoteFont\fontsize{7.5}{9}\selectfont,text=FigureMuted}]
% (a) chordal graph
\begin{scope}
 \node[v] (h1) at (0,9) {$1$}; \node[v] (h2) at (10,18) {$2$}; \node[v] (h3) at (10,0) {$3$};
 \node[v] (h4) at (22,18) {$4$}; \node[v] (h5) at (22,0) {$5$};
 \draw[e] (h1)--(h2) (h1)--(h3) (h2)--(h3) (h2)--(h4) (h3)--(h4) (h3)--(h5) (h4)--(h5);
 \node[sc] at (11,-10) {(a) 弦图};
\end{scope}
% (b) candidate graph with maximum-weight tree
\begin{scope}[xshift=40mm]
 \node[cl] (A) at (0,0) {$A=\{1,2,3\}$};
 \node[cl] (B) at (14,18) {$B=\{2,3,4\}$};
 \node[cl] (C) at (28,0) {$C=\{3,4,5\}$};
 \draw[FigureMuted,line width=.7pt,dashed] (A)--(C) node[midway,w]{$1$};
 \draw[FigureBlue,line width=1.6pt] (A)--(B) node[midway,w]{$2$};
 \draw[FigureBlue,line width=1.6pt] (B)--(C) node[midway,w]{$2$};
 \node[sc] at (14,-10) {(b) 最大权生成树，总权$4$};
\end{scope}
% (c) a wrong spanning tree
\begin{scope}[xshift=86mm]
 \node[cl] (B2) at (0,21) {$B=\{2,3,4\}$};
 \node[badcl] (A2) at (0,10.5) {$A=\{1,2,3\}$};
 \node[cl] (C2) at (0,0) {$C=\{3,4,5\}$};
 \draw[FigureStroke,line width=1pt] (B2)--(A2) (A2)--(C2);
 \node[sc] at (0,-10) {(c) 总权$3$};
 \node[note,align=center] at (0,-15) {含$4$的$B,C$被隔开};
\end{scope}
\end{tikzpicture}
